\documentclass[11pt]{article}
\usepackage[margin=2.4cm]{geometry}
\usepackage{amsmath,amssymb}
\usepackage{xcolor}
\usepackage{enumitem}
\usepackage{hyperref}

\newcommand{\comment}[1]{\medskip\noindent\textbf{#1}\par\nopagebreak\smallskip}
\newcommand{\reply}{\noindent\textit{Response.}\ }

\title{Response to Reviewers\\[2mm]
\large Manuscript: Graph-Based Reinforcement Learning for Electric Truck Fleet
Routing and Charging under Operational Uncertainty}
\date{}

\begin{document}
\maketitle

\noindent We thank the editor and both reviewers for reviews that were, in
several places, more accurate about our manuscript than we were. Three of the
criticisms turned out to identify real defects rather than presentational
weaknesses, and addressing them changed our own conclusions. We state those
three at the outset, because they alter claims that appeared in the original
submission and we would rather they be visible than discovered.

\begin{enumerate}[leftmargin=*]
\item \textbf{The feasibility mask does not explain our results.} Prompted by
Reviewer~1's comment~2, we implemented a training path without the hard mask.
Across three training seeds per arm, the masked and unmasked variants reach mean
success rates of $0.760$ and $0.773$ with overlapping ranges. The mask is a
convenience, not the source of the reported behaviour. What is necessary is that
infeasibility carry a cost commensurate with failure. The revised manuscript
says so in Section~5.6 and no longer attributes performance to masking.

\item \textbf{Our optimization benchmark was defective, and is now stronger.}
Prompted by Reviewer~1's comment~4 and Reviewer~2's comment~7, we validated the
nominal planner against exhaustive enumeration. On tiny deterministic instances
brute force beat our ``optimal'' planner on six of six scenarios, because its
circuit formulation could not leave a truck idle and therefore forced every
vehicle to serve at least one customer. After correction the planner matches
exhaustive enumeration on $30$ of $30$ instances, and every comparison in the
revised manuscript is against the corrected, stronger baseline.

\item \textbf{Training-seed variance exceeds most of the effects we previously
reported.} Replicating the identical configuration across seeds gives a spread
of $0.107$ in success rate. Several single-seed comparisons in the original
submission are inside that band and are no longer presented as findings. All
learned results now report either a seed range or an explicit single-seed
marker.
\end{enumerate}

\noindent Every experimental claim below is reproducible from artifacts
published with the revision, and each response states what changed, where, on
what evidence, and what remains open.

\section*{Handling Editor}

\comment{E1. The paper's claim that eVRP addresses single vehicles while eTFRP
addresses fleets is incorrect; there is extensive fleet-level eVRP literature,
including RL methods.}
\reply We accept this without reservation; the dichotomy was wrong as written.
The introduction no longer contrasts eTFRP with eVRP on the basis of fleet size.
The distinguishing features of our setting are now stated as the specific
combination we model, namely finite-port charging stations with endogenous
queueing, a nonlinear charging function integrated along the state-of-charge
trajectory, correlated travel and energy uncertainty, and event-driven
centralized control. We no longer claim that the existence of multiple vehicles
or shared chargers is itself novel.

\comment{E2. Operational uncertainty is well studied and the paper presents it
as an unfilled gap.}
\reply Accepted. The revision separates exogenous randomness (travel time,
energy draw, service time) from endogenous delay (queueing, which arises from
the fleet's own prior decisions), states the clipping bounds numerically, and
presents the contribution as the combination rather than as the treatment of
uncertainty per se. Table~2 now records for each quantity whether it is drawn
exogenously or arises endogenously.

We also report a limitation that this analysis exposed. In the evaluated
instances two trucks are distributed over twenty-five stations, so charger
contention almost never binds: an ablation that blanks the entire
charging-queue state leaves performance statistically unchanged. The endogenous
queueing mechanism is modelled and implemented, but the benchmark as configured
does not exercise it, and we now say so rather than claim credit for it.

\comment{E3. The methodological novelty is unclear; component-level evidence is
needed.}
\reply The contributions are restated as falsifiable component claims, each
tested at matched budget, observation width, curriculum, reward, and seed
stream, and re-scored on validation scenarios disjoint from those used for
checkpoint selection. Results appear in Section~5.6. Because independently
seeded runs of the identical configuration span $0.107$ in success rate, we
treat effects below roughly $0.1$ as indistinguishable from training noise and
report them as such.

Three components clear that threshold: the per-action routing features
($0.213$ when removed, against a $0.700$--$0.807$ seed band), the typed pairwise
relations ($0.633$), and the pooled state embedding ($0.667$). Two do not: the
charging-queue state ($0.793$) and the active-truck indicator ($0.727$). We
report the negative results alongside the positive ones. On scalability, we no
longer claim a computational advantage in the joint setting: the nominal model
is small enough that CP-SAT proves optimality in milliseconds, and a budget
sweep from $0.5$~s to $45$~s returns an identical plan.

\section*{Reviewer 1}

\comment{R1.1. In the main eTFRP experiments, deliveries are pre-assigned and
mostly ordered, and depot return is not required, so the problem is closer to
energy-feasible execution than to fleet routing.}
\reply The observation is correct. We have kept the eTFRP as the principal
setting because it reflects the operational case this work targets, in which
assignments follow from contracts and depot logistics rather than from the
routing system, and we now say this explicitly rather than implying that the
harder combinatorial problem is being solved.

To address the substance of the criticism, we added a second setting in which
precisely the removed decisions are restored: customers are owned by the fleet,
assignment and sequencing are decided online, payload capacity binds, and every
truck must return to the depot. It replaces the former single-truck eVRP study
as the final experimental section, since it answers this comment directly
whereas the single-vehicle study did not. The same architecture is applied
unchanged and reaches $84.3\%$ success $[83.0, 85.8]$ over three seeds against
$62.2\%$ for the corrected CP-SAT planner, while differing from that planner by
$+0.2$ hours $[-3.1, +3.7]$ in paired fleet travel time on jointly solved
scenarios. See Section~5.5 and Table~7.

\comment{R1.2. The feasibility mask may explain most of GraphPPO's gain.}
\reply This was the most consequential comment in the review, and the answer is
not the one we expected. We implemented a mask-free training path in which the
observation and candidate set are identical and only slots denoting no action at
all are hidden, so the policy must learn feasibility from experience.

Across three seeds per arm the masked and unmasked variants are
indistinguishable: $0.760$ against $0.773$ mean success, with per-arm ranges of
$0.107$ and $0.113$. The unmasked policy averages $0.1$ invalid actions per
episode, i.e.\ it learns feasibility on its own. What proves necessary is the
\emph{cost} of infeasibility: when an infeasible action is refused at a tenth of
the failure penalty the policy never learns feasibility and success collapses to
zero, while at the full failure penalty it recovers to $0.773$.

We also compare, under equal information and identical budgets, a flat-state
MaskPPO analogue, DeepSets-PPO, a state-GNN with independent action scoring, and
a constructive attention model. The action head, not the state encoder, is what
separates the family: every independent-head arm lands at $0.527$--$0.573$,
while the attention encoder under a complete-GCN head reaches $0.773$, inside
the graph encoder's seed band. We therefore no longer attribute our results to
the graph state encoder either. Reported metrics are feasibility, conditional
travel time and makespan, queue time, and per-decision runtime, never reward
alone.

\comment{R1.3. Notation and equation defects.}
\reply All corrected, and the equations were regenerated from the tested
implementation rather than edited in place.
\begin{itemize}[leftmargin=*,nosep]
\item Nominal and realized quantities now use distinct symbols consistently:
$\tau_{uv}, b_{uv}$ against $\tilde{\tau}_{uv}, \tilde{b}_{uv}$.
\item Clipping bounds are given numerically:
$\tilde{b}_{uv}\in[0.90 b_{uv}, 1.20 b_{uv}]$ and
$\tilde{\tau}_{uv}\in[0.85\tau_{uv}, 2.50\tau_{uv}]$, applied to the realized
quantity only, so the nominal network read by the planners is unaffected.
\item The charging constraint is rewritten with the state-of-charge trajectory
$\mathrm{soc}_i(t+s)$ and elapsed-time integration variable $s$. We note why it
matters: evaluating the power curve at the initial state overstates delivered
energy by up to $55\%$ for a session crossing the taper.
\item Equation~19 used the receiver embedding twice. Messages now condition on
the sender embedding and the edge features, with the receiver entering at the
update step.
\item Equation~22 summed the node's own embedding over its neighbors. It now
aggregates neighboring action embeddings, with self-influence carried by a
separate term and the single-feasible-action case defined explicitly.
\end{itemize}

\comment{R1.4. The optimization benchmark does not establish near-optimality.}
\reply Accepted, and the investigation prompted by this comment found a defect
in our own baseline, reported in the preamble. The inherited per-truck models
are now named for what they are in the source and are no longer presented as
optimality references. The campaign baseline is a fleet-level CP-SAT model whose
objective matches exhaustive enumeration on $30$ of $30$ tiny instances; an
adaptive large neighbourhood search is included as an independent check and
reaches the same optimum on all of them. Every episode now publishes solver
status, objective, best bound, relative gap, wall time, solve count, and the
number of times the executed plan fell back to the shared navigation layer.

We have removed the term ``near-optimal'' where it was unsupported. A bounded
optimality statement is available only on instances small enough to enumerate;
at campaign scale we report the best-known plan and say so.

\comment{R1.5. High reward is reported despite zero success rate.}
\reply Accepted. Reward is no longer a headline metric anywhere. Evaluation is
feasibility-first: success rate with Wilson intervals dominates, and cost
metrics are reported with their conditioning stated. Every failed episode is
retained in every non-conditioned aggregate and is now labelled with a cause;
truncation was previously silent, which is where unlabelled failures came from.
The inherited shaping has been replaced and every coefficient is published.

\comment{R1.6. Baselines are weak and underspecified.}
\reply Added: an ALNS metaheuristic with four destroy operators, greedy and
regret-2 repair, adaptive operator weights and simulated-annealing acceptance,
searching the same nominal costs as CP-SAT and executing through the same
energy-safe layer; a constructive attention baseline; DeepSets-PPO; and a
state-GNN with independent action scoring. All share the observation, mask,
curriculum, reward, seed stream and interaction budget. The claim that PPO is
``the most capable discrete-action RL algorithm'' has been removed.

Two honest notes. ALNS matches CP-SAT closely, which we take as evidence that
the nominal problem is solved rather than that either method is weak. And the
attention baseline is competitive with our own encoder, which is why we no
longer claim the graph encoder as a contribution.

\comment{R1.7. Generalization evidence is too narrow.}
\reply The former ``zero-shot'' study is renamed size transfer, because that is
what it measures: only instance size varies. Genuine distribution shift is now
studied separately across twenty regimes, each scoring every method on identical
held-out scenarios and each labelled interpolation, size transfer, or
out-of-distribution so that one kind of evidence is never quoted as another. The
main setting is evaluated on $500$ paired test scenarios rather than $300$, and
learned results carry seed ranges.

The finding is bounded rather than uniformly favourable. Transfer holds across
instance size and across shifts in charger power, efficiency, port availability,
travel and energy variance, demand, and road-network speed, with our method
leading thirteen of sixteen regimes. It fails where the resource budget itself
moves: with a $300$~kWh battery we reach $0.34$ against CP-SAT's $0.43$, with
$500$~kWh $0.87$ against $0.98$, and with doubled service times $0.46$ against
$0.60$. Averaged over out-of-distribution regimes our success changes by
$-0.110$ against $-0.006$ for ALNS. We state this in the conclusion as a
deployment boundary.

\section*{Reviewer 2}

\comment{R2.1--R2.3. The problem definition is confusing; uncertainty and
mandatory service should be stated early.}
\reply The introduction now defines the joint problem in operational terms
without the eTFRP/eVRP dichotomy, and enumerates the stochastic and endogenous
quantities where the problem is introduced rather than later. Every-customer-once
service, payload capacity, energy feasibility, the optional time-window variant,
and depot return are stated as explicit constraints, and the outcome taxonomy
(success, infeasibility, timeout, incomplete service) is defined.

\comment{R2.4--R2.6. Language, reward coefficients, charging discretization.}
\reply The phrase ``stochastic edge traversing'' is replaced by ``realized
travel duration on the selected road-network transition''. Every training reward
coefficient is published. On discretization, we now have evidence rather than a
choice: refining the target state-of-charge grid from $10\%$ to $5\%$ leaves
performance inside the seed band, whereas expressing charging as a $15/30/60$
minute duration degrades it to $0.627$ success and $159.3$ travel hours. A fixed
duration delivers different energy depending on position on the charging curve,
so the same action denotes different outcomes in different states; a target
state of charge does not.

\comment{R2.7. The optimization model is unclear.}
\reply A complete formulation is given with its queueing, uncertainty, horizon,
information and fallback assumptions stated, together with which baselines are
offline, rolling-horizon, deterministic or scenario-based. Tiny exact objectives
are validated against exhaustive enumeration, which is how we found the defect
described in the preamble.

\comment{R2.8--R2.9. Heuristic and neural baselines; attention and transformer
methods should be compared.}
\reply A constructive attention baseline is included, adapted to this setting
and with the adaptation stated rather than glossed: Kool et al.\ assume a single
vehicle, no charging, and no exogenous uncertainty, so a literal port is
impossible; what transfers is the architecture, and we give the typed edge
features to it as an attention bias so that it reads exactly the same content as
our encoder. Each implemented heuristic is linked to the closest cited variant
with its adaptations disclosed.

\comment{Minor comments.}
\reply ``demonstrated'' is now ``demonstrates''; the missing comma in the
objective equation is added; figure caption capitalization is standardized;
Table~2 reports realized edge energy rather than only the coefficient $\xi$. The
panel label ``c.\ Actor Network head'' is in the figure source asset and is
corrected to ``(c) Actor Network Head'' in the regenerated figure.

\section*{Additional comments}

\comment{Unloading/service time as a model parameter; endogenous queue delay;
sourcing and validating the charging equation; comparison with a three-segment
piecewise formulation; an expanded random-variable table.}
\reply Service time is defined as a model parameter, sampled once per service
event, with its distribution and cap in Table~2. Queue delay is explained as
endogenous under finite ports and FCFS admission, with the honest caveat under
E2 that the evaluated fleet-to-station ratio does not exercise it.

On the charging function we now report a quantitative comparison rather than an
assertion. Against the integrated CCCV curve at $350$--$750$~kW, a constant-power
model errs by $5.0\%$ on average and $28.7\%$ at worst; the classical
three-segment piecewise-linear formulation of the Montoya type errs by $4.7\%$
and $27.8\%$, i.e.\ it is barely better. The reason is structural and we state
it: our curve ramps to peak power before tapering and is therefore not concave,
while a Montoya-style approximation assumes concavity. Placing breakpoints
inside both curved regions reduces the error to $0.6\%$ mean and $4.0\%$
maximum. We also fixed a defect found during this comparison, in which the
charge-time estimator failed to converge for a full charge and returned $10$
hours for a session requiring $0.53$ hours.

Table~2 is expanded to give, for each random variable, its distribution,
parameters, clipping bounds, correlation structure, and whether it is exogenous
or endogenous.

\section*{Remaining limitations}

We prefer to name these rather than leave them to be found.
\begin{itemize}[leftmargin=*,nosep]
\item Charger contention is modelled but not exercised by the evaluated
instances, so no claim about queue-aware behaviour is supported by our results.
\item Optimality is bounded only on instances small enough to enumerate.
\item Out-of-distribution performance degrades where the energy or time budget
moves, and in those regimes a re-solving planner is the better choice.
\item Learned results use three training seeds. Given the measured variance we
consider three sufficient for the final configuration, whose spread is $0.026$,
but not for the intermediate training stages, whose spread is $0.120$ and whose
rankings we therefore no longer interpret.
\end{itemize}

\end{document}
